\color{unimain}{\subsection{Логика отключения выключателей}}
\color{black}

\begin{enumerate}

\item{Общие сведения}

\begin{enumerate}

\item
Ввод/вывод функции осуществляется либо посредством программного ключа \textbf{«ЛО Вn»}, либо с использованием внешнего сигнала \textit{«Вывод действия на Вn»}.

\item
Уставки логики отключения В1 представлены в таблице \ref{app_set:lov1}. Уставки логики отключения В1 представлены в таблице \ref{app_set:lov2}.

\end{enumerate}


%%%%%%%%%%%%%%%%%%%%%%%%%%%%%%%%%%%%%%%%%%%%%%%%%% Логика отключения %%%%%%%%%%%%%%%%%%%%%%%%%%%%%%%%%%%%%%%%%%%%%%%%%%
\item{Логика отключения}

\begin{enumerate}

\item
Функциональный блок логики отключения предназначен для подачи сигнала на отключение силового выключателя.

\item
Команда на отключение выключателя \textit{«Отключить Вn»} формируется:
\begin{itemize}
\item при отключении ошиновки (\textit{«Отключение СШ»});
\item при отключении в результате опробования присоединения;
\item от команды оперативного отключения выключателя (\textit{«РКО Вn»});
\item при действии УРОВ на отключение «своего» выключателя (\textit{«УРОВ Вn на себя»}).
\end{itemize}

\item
При необходимости может быть введен подхват команды на отключение от токовых реле, установленных в цепях ЭМО, в этом случае сигнал на отключение удерживается в сработанном состоянии до
исчезновения тока через ЭМО (пока блок-контакты выключателя не разорвут цепи отключения). Ввод осуществляется программным ключом \textbf{«Подхват откл. Вn от РТ ЭМО»}. 

\item
При наличии сигнала \textit{«Блок. отключения Вn»}, формируемого функциональным блоком контроля силового выключателя отключение выключателя блокируется.

\item
Выдержка времени на возврат \textbf{«Тпрод откл. Вn»} задает время продления сигнала отключения.

\item
Алгоритм работы логики отключения В1 приведена на рисунке \ref{pict_lo:LO_Q1}, алгоритм работы логики отключения В2 выполнен аналогично. 

\medskip
\begin{figure}[H]
\centering
\includegraphics[width=0.8\textwidth,height=0.8\textheight,keepaspectratio]{\fbpath/05.02 ДЗО/761.9901 Логика отключения В/img/LO_Q1.pdf}
\captionof{figure}{Функциональная схема логики отключения В1}\label{pict_lo:LO_Q1}
\end{figure}


\end{enumerate}


%%%%%%%%%%%%%%%%%%%%%%%%%%%%%%%%%%%%%%%%%%%%%%%%%% Логика запрета АПВ %%%%%%%%%%%%%%%%%%%%%%%%%%%%%%%%%%%%%%%%%%%%%%%%%%
\item{Логика запрета АПВ}

\begin{enumerate}

\item
Сигналы запрета АПВ формируется в следующих условиях:
\begin{itemize}
\item от сигнала об аварийном снижении давления элегаза в ТТ;
\item при отключении при опробовании;
\item от общего сигнала запрета АПВ;
\item от команды оперативного отключения выключателя;
\item при срабатывании ЗНФ (при введенной уставку \textbf{«Запрет АПВ Вn от ЗНФ Вn»});
\item если выключатель не включился за время \textbf{«Tожид. усл. пуска АПВ Вn»} (при введенном программном ключе \textbf{«Огр. ожид. пуска АПВ Вn»});
\item если смежный выключатель (ведущий) включился на КЗ.
\end{itemize}

\item
Алгоритм работы логики запрета АПВ В1 приведена на рисунке \ref{pict_lo:LZ_APV}, алгоритм работы логики запрета АПВ В2 выполнен аналогично. 

\medskip
\begin{figure}[H]
\centering
\includegraphics[width=0.8\textwidth,height=0.8\textheight,keepaspectratio]{\fbpath/05.02 ДЗО/761.9901 Логика отключения В/img/LZ_APV.pdf}
\captionof{figure}{Функциональная схема логики запрета АПВ В1}\label{pict_lo:LZ_APV}
\end{figure}


\end{enumerate}


\end{enumerate}